% ============================================================
%  preamble.tex —— 「李思风」极简学术笔记样式
%  参照《线性代数：原理与应用》（李思，清华大学）讲义排版：
%    · 正文方正宋体 11pt、1.3 倍行距、首行缩进 2 字符
%    · 章标题居中加粗特大、节标题居中、小节左对齐，纯黑无装饰
%    · 定义/定理/命题/引理/推论 = 圆角细边框白底盒（行内加粗彩色标签）
%    · 例/练习/注/证明 = 无框，黑色加粗行内标签（注用楷体）
%    · 页码底部居中、无页眉；目录默认点线引导 + 蓝色超链接
%  编译方式：XeLaTeX（Fandol 字体随 TeX Live 提供，跨平台一致）
% ============================================================

% ------------------------------------------------------------%
% 基础：字体、页面
% ------------------------------------------------------------%
\usepackage[fontset=fandol]{ctex}   % 方正宋体/楷体/黑体，与参照讲义同款
\usepackage{geometry,graphicx,xcolor,color,subcaption}
\geometry{a4paper, margin=2.54cm, headheight=14.5pt}

% ------------------------------------------------------------%
% 数学与表格
% ------------------------------------------------------------%
\usepackage{amssymb,amsmath,amsthm}
\usepackage{cancel}
\usepackage{booktabs}
\usepackage{multirow}
\usepackage{newclude,ulem}

% ------------------------------------------------------------%
% 配色：绿=定义，蓝=定理，深蓝=命题/引理，橙=推论（李思讲义取色）
% ------------------------------------------------------------%
\definecolor{LiSiDef}{RGB}{0,191,0}       % 定义：绿
\definecolor{LiSiThm}{RGB}{0,0,191}       % 定理：蓝
\definecolor{LiSiPro}{RGB}{0,0,153}       % 命题/引理：深蓝
\definecolor{LiSiCor}{RGB}{191,96,0}      % 推论：橙
\definecolor{LiSiWarn}{RGB}{191,0,0}      % 注意：红

% 兼容旧配色名（艺术封面与历史章节文件仍引用）
\definecolor{winered}{rgb}{0.5,0,0}
\definecolor{structurecolor}{RGB}{24,44,82}
\definecolor{gold}{RGB}{168,132,74}
\definecolor{main}{rgb}{0.5,0,0}
\definecolor{second}{RGB}{115,45,2}
\definecolor{third}{RGB}{0,80,80}
\colorlet{thmColor}{LiSiThm}
\colorlet{defColor}{LiSiDef}
\colorlet{proColor}{LiSiPro}
\colorlet{exColor}{LiSiPro}
\colorlet{keyColor}{LiSiPro}
\colorlet{tipColor}{LiSiDef}
\colorlet{warnColor}{LiSiWarn}
\colorlet{thmBg}{white}
\colorlet{defBg}{white}
\colorlet{exBg}{white}
\colorlet{proBg}{white}
\colorlet{keyBg}{white}
\colorlet{tipBg}{white}
\colorlet{warnBg}{white}

\usepackage[colorlinks,linkcolor=blue,citecolor=blue,urlcolor=blue]{hyperref}

% ------------------------------------------------------------%
% 计数器：定理类（定义/定理/命题/引理/推论）共用；例/练习独立
% 章号为中文「第一章」仅用于显示，编号仍为阿拉伯数字（1.1、2.3.1）
% ------------------------------------------------------------%
\makeatletter
\@ifundefined{c@chapter}{%
    % ---------- article 类：无章，节为最高级 ----------
    \newcounter{theorem}%
    \newcounter{example}%
}{%
    % ---------- book 类：章为最高级 ----------
    \renewcommand{\thechapter}{第\chinese{chapter}章}       % 章号：第一章、第二章（用于标题、目录与引用）
    \renewcommand{\thesection}{\arabic{chapter}.\arabic{section}}
    \renewcommand{\thesubsection}{\thesection.\arabic{subsection}}
    \renewcommand{\thesubsubsection}{\thesubsection.\arabic{subsubsection}}
    \renewcommand{\theequation}{\arabic{chapter}.\arabic{equation}}
    \renewcommand{\thefigure}{\arabic{chapter}.\arabic{figure}}
    \renewcommand{\thetable}{\arabic{chapter}.\arabic{table}}
    \newcounter{theorem}[chapter]
    \renewcommand{\thetheorem}{\arabic{chapter}.\arabic{theorem}}
    \newcounter{example}[chapter]
    \renewcommand{\theexample}{\arabic{chapter}.\arabic{example}}
}
\makeatother

% ------------------------------------------------------------%
% 定理类环境：圆角细边框白底盒 + 行内加粗彩色标签
% ------------------------------------------------------------%
\usepackage{etoolbox}
\usepackage{tcolorbox}
\tcbuselibrary{skins,breakable}

\tcbset{
    lisiBox/.style={
        enhanced, breakable,
        colback=white, boxrule=1.4pt, arc=4pt,
        left=4mm, right=4mm, top=2mm, bottom=2mm,
        before skip=1.3em, after skip=1.3em,
        fontupper=\linespread{1.4}\selectfont
    }
}

% 生成带编号的定理类环境：#1 环境名，#2 标签中文，#3 边框/标签色
\newcommand{\newlisiThm}[3]{%
    \newenvironment{#1}[1][]{%
        \refstepcounter{theorem}\par
        \begin{tcolorbox}[lisiBox, colframe=#3]%
        \noindent{\bfseries\songti\color{#3}#2\ \thetheorem\ifstrempty{##1}{}{（##1）}}%
        \hspace{0.4em}\ignorespaces
    }{\end{tcolorbox}}%
}
\newlisiThm{definition}{定义}{LiSiDef}
\newlisiThm{axiom}{公理}{LiSiDef}
\newlisiThm{theorem}{定理}{LiSiThm}
\newlisiThm{proposition}{命题}{LiSiPro}
\newlisiThm{lemma}{引理}{LiSiPro}
\newlisiThm{property}{性质}{LiSiPro}
\newlisiThm{corollary}{推论}{LiSiCor}

% 例题 / 练习：无框，黑色加粗行内标签（例 x.y.z.）
\newenvironment{example}[1][]{%
    \par\vspace{0.5em}\refstepcounter{example}%
    \noindent{\bfseries 例\ \theexample.\ifstrempty{#1}{}{\ #1}}\hspace{0.4em}\ignorespaces
}{\par\vspace{0.2em}}
\newenvironment{longexample}[1][]{%
    \par\vspace{0.5em}\refstepcounter{example}%
    \noindent{\bfseries 例\ \theexample.\ifstrempty{#1}{}{\ #1}}\hspace{0.4em}\ignorespaces
}{\par\vspace{0.2em}}
\newenvironment{exercise}[1][]{%
    \par\vspace{0.5em}\refstepcounter{example}%
    \noindent{\bfseries 练习\ \theexample.\ifstrempty{#1}{}{\ #1}}\hspace{0.4em}\ignorespaces
}{\par\vspace{0.2em}}

% 无编号盒子：重点 / 技巧 / 注意（细边框 + 行内加粗标签）
\newtcolorbox{keybox}[1][重点]{lisiBox, colframe=LiSiPro,
    before upper={\noindent{\bfseries\songti\color{LiSiPro}#1}\hspace{0.4em}}}
\newtcolorbox{tipbox}[1][技巧]{lisiBox, colframe=LiSiDef,
    before upper={\noindent{\bfseries\songti\color{LiSiDef}#1}\hspace{0.4em}}}
\newtcolorbox{warnbox}[1][注意]{lisiBox, colframe=LiSiWarn,
    before upper={\noindent{\bfseries\songti\color{LiSiWarn}#1}\hspace{0.4em}}}
\newtcolorbox{myTheorem}[1]{lisiBox, colframe=LiSiPro,
    before upper={\noindent{\bfseries\songti\color{LiSiPro}#1}\hspace{0.4em}}}
\newtcolorbox{mydefinition}[1]{lisiBox, colframe=LiSiDef,
    before upper={\noindent{\bfseries\songti\color{LiSiDef}#1}\hspace{0.4em}}}
\newtcolorbox{myhypothesis}[1]{lisiBox, colframe=LiSiCor,
    before upper={\noindent{\bfseries\songti\color{LiSiCor}#1}\hspace{0.4em}}}

% 公式强调框（无标题，内容居中）
\newtcolorbox{myformula}[1][]{
    enhanced, breakable, colback=white, boxrule=1pt, arc=4pt,
    colframe=black!35,
    before skip=1.2em, after skip=1.2em,
    center upper,
    top=2.4mm, bottom=2.4mm, left=4mm, right=4mm,
    #1
}

% ------------------------------------------------------------%
% 证明 / 解 / 注 / 引言落款
% ------------------------------------------------------------%
\renewenvironment{proof}[1][证明]{%
    \par\vspace{0.5em}\noindent{\bfseries\songti #1：}\hspace{0.3em}\ignorespaces
}{\hfill$\square$\par\vspace{0.3em}}
\newenvironment{solution}[1][解]{%
    \par\vspace{0.5em}\noindent{\bfseries\songti #1．}\hspace{0.3em}\ignorespaces
}{\par\vspace{0.3em}}
\newenvironment{remark}[1][注]{%
    \par\vspace{0.5em}\noindent{\bfseries\songti #1．}\hspace{0.3em}\kaishu\ignorespaces
}{\par\vspace{0.3em}}
\newenvironment{note}[1][注]{%
    \par\vspace{0.5em}\noindent{\bfseries\songti #1：}\hspace{0.3em}\kaishu\ignorespaces
}{\par\vspace{0.3em}}
\newcommand{\intro}[1]{\rightline{\parbox[t]{5cm}{\footnotesize\fangsong\quad\quad #1}}}

% ------------------------------------------------------------%
% 列表：默认黑色编号/圆点，紧凑间距
% ------------------------------------------------------------%
\usepackage{enumerate}
\usepackage{enumitem}
\setlist{topsep=0.35em, itemsep=0.2em, parsep=0.15em, partopsep=0.1em}
\setlist[enumerate,1]{label=\arabic*.}
\setlist[enumerate,2]{label=(\arabic*)}
\setlist[enumerate,3]{label=\Roman*.}
\setlist[enumerate,4]{label=\Alph*.}

\usepackage{float}
\usepackage{tikz}
\usetikzlibrary{angles, quotes, arrows.meta, calc}

% ------------------------------------------------------------%
% 代码（Python）
% ------------------------------------------------------------%
\usepackage{listings}
\lstdefinestyle{pythonstyle}{
    language=Python,
    basicstyle=\ttfamily\small,
    keywordstyle=\color{blue}\bfseries,
    commentstyle=\color{green!50!black}\itshape,
    stringstyle=\color{red},
    numberstyle=\tiny\color{gray},
    numbers=left,
    numbersep=8pt,
    frame=single,
    rulecolor=\color{black!30},
    breaklines=true,
    showstringspaces=false,
    tabsize=4,
    captionpos=b,
    xleftmargin=2em,
    framexleftmargin=1.5em
}
\lstset{style=pythonstyle}

% ------------------------------------------------------------%
% 标题格式：章居中特大加粗、节居中、小节左对齐，纯黑无装饰
% ------------------------------------------------------------%
\usepackage{titlesec}
\usepackage{titletoc}
\linespread{1.3}

% 目录：章条目「第一章 标题」（章号不做定宽盒，避免中文标签重叠）
\makeatletter
\@ifundefined{c@chapter}{}{%
    \titlecontents{chapter}[0em]{\addvspace{1.1em}\bfseries}
        {\thecontentslabel\hspace{1em}}{}{\hfill\contentspage}

    \titleformat{\chapter}[block]
        {\normalfont\centering\bfseries\songti\fontsize{20.7}{26}\selectfont}
        {\thechapter}{1em}{}
    \titlespacing*{\chapter}{0pt}{52pt plus 4pt minus 4pt}{40pt plus 4pt minus 4pt}%
}
\makeatother

\titleformat{\section}[block]
    {\normalfont\centering\bfseries\songti\Large}
    {\thesection}{1em}{}
\titlespacing*{\section}{0pt}{18pt plus 3pt minus 3pt}{10pt plus 2pt minus 2pt}

\titleformat{\subsection}[hang]
    {\normalfont\bfseries\songti\large}
    {\thesubsection}{1em}{}
\titlespacing*{\subsection}{0pt}{14pt plus 2pt minus 2pt}{6pt plus 2pt minus 2pt}

\titleformat{\subsubsection}[hang]
    {\normalfont\bfseries\songti\normalsize}
    {\thesubsubsection}{1em}{}
\titlespacing*{\subsubsection}{0pt}{10pt plus 2pt minus 2pt}{4pt plus 1pt minus 1pt}

% ------------------------------------------------------------%
% 页眉页脚：无页眉，页码底部居中（plain / mainstyle 同为该样式）
% ------------------------------------------------------------%
\usepackage{fancyhdr}
\fancypagestyle{plain}{%
    \fancyhf{}
    \fancyfoot[C]{\normalsize\thepage}
    \renewcommand{\headrulewidth}{0pt}
    \renewcommand{\footrulewidth}{0pt}
}
\fancypagestyle{mainstyle}{%
    \fancyhf{}
    \fancyfoot[C]{\normalsize\thepage}
    \renewcommand{\headrulewidth}{0pt}
    \renewcommand{\footrulewidth}{0pt}
}
\pagestyle{plain}

% ------------------------------------------------------------%
% 封面主题色（默认封面 cover.tex 使用；其余封面在各自文件内定义）
% ------------------------------------------------------------%
\definecolor{DarkBlue}{RGB}{0,0,139}

% --- 「印象·日出」封面配色（cover_Impression风.tex 使用，莫奈《日出·印象》取色） ---
\definecolor{imSky}{RGB}{247,242,232}       % 晨空米白（画面基底）
\definecolor{imMist}{RGB}{210,203,192}      % 远景晨雾灰
\definecolor{imSun}{RGB}{224,110,84}        % 日出橙红（太阳主体）
\definecolor{imSunGlow}{RGB}{238,160,120}   % 太阳柔光晕
\definecolor{imSea}{RGB}{74,104,128}        % 海面青灰蓝（水面主色）
\definecolor{imSeaDark}{RGB}{52,78,100}     % 海面深蓝灰（波纹暗部/剪影）
\definecolor{imSeaLight}{RGB}{142,168,186}  % 海面浅蓝灰（波纹亮部）

% --- 「格尔尼卡」封面配色（cover_Guernica风.tex 使用，毕加索《格尔尼卡》取色：纯灰阶） ---
\definecolor{gcPaper}{RGB}{235,232,226}     % 灰白画布（画面基底）
\definecolor{gcPaperDark}{RGB}{206,202,194} % 画布暗部（下缘）
\definecolor{gcBone}{RGB}{247,245,241}      % 骨白（受光碎片/提亮）
\definecolor{gcLight}{RGB}{192,188,180}     % 浅灰（中间调碎片）
\definecolor{gcMid}{RGB}{130,128,124}       % 中灰（暗部碎片）
\definecolor{gcChar}{RGB}{62,62,64}         % 炭灰（剪影/线条）
\definecolor{gcBlack}{RGB}{24,24,26}        % 墨黑（重点剪影/标题）
